\chapter{Fase 2: especificação da avaliação}
\label{cap:fase2}

A Fase 2 especifica a avaliação: seleciona as métricas e estabelece os níveis de pontuação e os critérios de julgamento \cite{iso25040,ramos2026processo}. A especificação usa o método GQM \cite{basili1994} e foi escrita para que a avaliação seja repetível: cada métrica tem definição, cálculo, fonte dos dados, procedimento de coleta e interpretação documentados na seção ``Fase 2'' da GitPage. Este capítulo apresenta a síntese dessas decisões.

\section{Rastreabilidade com a Fase 1}
\label{sec:rastreabilidade}

Todos os elementos do GQM derivam de decisões da Fase 1 (\autoref{tab:rastreabilidade}). As quatro características priorizadas originam os objetivos; o ponto de vista e os perfis definem para quem as métricas são calculadas; as subcaracterísticas mantidas originam as questões; e as observações preliminares fundamentam as hipóteses.

\begin{table}[htb]
	\caption{Rastreabilidade entre as decisões da Fase 1 e a especificação da Fase 2}
	\label{tab:rastreabilidade}
	\centering
	\scriptsize
	\begin{tabularx}{\textwidth}{L{5.2cm} Y}
		\toprule
		\textbf{Decisão da Fase 1} & \textbf{Uso na Fase 2} \\
		\midrule
		Características priorizadas (\autoref{tab:matriz}) & Um objetivo por característica (G1 a G4) \\
		Subcaracterísticas mantidas no modelo adaptado & Uma ou mais questões por subcaracterística (Q1 a Q9, Q11 a Q14) \\
		Ponto de vista do estudante e perfis (calouro, veterano, formando, uso só pelo celular, com deficiência) & Ponto de vista dos objetivos; resultados por perfil em M5, M6, M7 e M11; inclusividade em Q8; celular em Q12; questionário por semestre \\
		Observações preliminares OP-01 a OP-08 (\autoref{tab:op}) & Fundamento das hipóteses: OP-01 em M11, OP-02 e OP-03 em M14, OP-04 em M4, OP-05 em M6 e M7, OP-06 em M8 e M20, OP-07 em M1 e M3 \\
		Decisões apoiadas (\autoref{tab:decisoes}) & Resultados por tarefa e ranking das tarefas; G3 apoia a decisão sobre interface responsiva; M16 a M19 apoiam a otimização de páginas \\
		Tarefas T1 a T6 e ambientes (\autoref{tab:tarefas}) & Tarefas usadas em todas as medições; navegadores, telas e rede em G3 e G4 \\
		Rigor intermediário & Dois avaliadores em M1, M2, M10 e M20; cinco execuções em M16 a M19; critérios fixados antes da coleta \\
		\bottomrule
	\end{tabularx}
	\fonte{Elaborado pela equipe.}
\end{table}

\section{Objetivos de medição}
\label{sec:objetivos-gqm}

Cada objetivo foi escrito no \textit{template} GQM (\autoref{tab:objetivos}). Em todos, o propósito é \textbf{avaliar, a fim de identificar pontos fracos}, e o ponto de vista é o do \textbf{estudante de graduação da UnB}, com o recorte de perfil indicado. A numeração não segue a prioridade da Fase 1, que é usada na apresentação dos resultados: G2, G1, G4 e G3.

\begin{table}[htb]
	\caption{Objetivos de medição no \textit{template} GQM}
	\label{tab:objetivos}
	\centering
	\scriptsize
	\begin{tabularx}{\textwidth}{L{1.7cm} Y Y Y Y}
		\toprule
		& \textbf{G1 -- Adequação Funcional} & \textbf{G2 -- Capacidade de Interação} & \textbf{G3 -- Flexibilidade} & \textbf{G4 -- Eficiência de Desempenho} \\
		\midrule
		\textbf{Analisar} & o portal do discente do SIGAA/UnB & a interface web do portal do discente & a interface web do portal do discente & as páginas usadas nas tarefas T1 a T6 \\
		\textbf{Com respeito à} & completude, correção e adequação (pertinência) funcional & reconhecimento de adequação, apreensibilidade, operabilidade, proteção contra erros, inclusividade, autodescritividade e satisfação & adaptabilidade a navegadores, telas e toque & comportamento temporal e utilização de recursos no cliente \\
		\textbf{Ponto de vista (recorte)} & estudante de graduação & estudante, com atenção a calouros e a pessoas com deficiência & estudante que acessa só pelo celular & estudante que acessa por rede móvel \\
		\textbf{Contexto} & tarefas T1 a T6 e planejamento do fluxo, versão v4.15.10.115, 2026/2 & execução de T1 a T6 em computador & Chrome, Firefox, Edge e Safari; computador, tablet e celular & medição no cliente, em rede comum e móvel simulada \\
		\bottomrule
	\end{tabularx}
	\fonte{Elaborado pela equipe.}
\end{table}

\section{Questões e hipóteses}
\label{sec:questoes}

Cada questão caracteriza um atributo do foco de qualidade do seu objetivo, e cada uma recebe uma hipótese com critério de confirmação numérico, declarado antes da coleta (\autoref{tab:questoes}). Na Fase 4, cada hipótese será classificada como confirmada, refutada ou inconclusiva; uma hipótese refutada não é falha da avaliação, mas indica que a percepção da Fase 1 não se confirmou.

{\scriptsize
\begin{longtable}{L{0.55cm} L{4.5cm} L{2.4cm} L{7.0cm}}
	\caption{Questões, atributos e hipóteses}
	\label{tab:questoes} \\
	\toprule
	\textbf{ID} & \textbf{Questão} & \textbf{Atributo} & \textbf{Hipótese (critério de confirmação)} \\
	\midrule
	\endfirsthead
	\toprule
	\textbf{ID} & \textbf{Questão} & \textbf{Atributo} & \textbf{Hipótese (critério de confirmação)} \\
	\midrule
	\endhead
	\bottomrule
	\endfoot
	\multicolumn{4}{l}{\textit{G1 -- Adequação Funcional}} \\
	Q1 & O SIGAA oferece as funções de que o estudante precisa? & Completude & H1: cobertura entre 75\% e 90\%, pior em T4 e no planejamento do fluxo (M1 Regular) \\
	Q2 & As informações exibidas estão corretas? & Correção & H2: nenhuma divergência com as fontes oficiais (M2 = 100\%) \\
	Q3 & As funções bastam para planejar o fluxo sem ferramentas externas? & Pertinência & H3: entre 25\% e 50\% dos respondentes usam ferramenta externa (M3 Regular) \\
	\multicolumn{4}{l}{\textit{G2 -- Capacidade de Interação}} \\
	Q4 & O estudante identifica onde está a função? & Reconhecimento de adequação & H4: entre 50\% e 75\% de acerto na primeira escolha de menu, pelo menu Ensino extenso (M4 Regular) \\
	Q5 & Conclui as tarefas sem ajuda, inclusive no primeiro contato? & Apreensibilidade & H5: conclusão geral entre 78\% e 90\% (M5 Bom), mas abaixo de 65\% entre calouros \\
	Q6 & A navegação evita esforço desnecessário? & Operabilidade & H6: tempo relativo entre 2,5 e 4,0 e desvio de caminho entre 1,33 e 2,00 (M6 e M7 Regular) \\
	Q7 & O sistema evita perda de trabalho após erros ou expiração da sessão? & Proteção contra erros & H7: no máximo dois de seis cenários tratados (M20 Insatisfatório) e recuperação de 50\% a 70\% dos erros (M8 Regular) \\
	Q8 & As páginas podem ser usadas por pessoas com deficiência? & Inclusividade & H8: mais de cinco erros automáticos por página (M9 Insatisfatório) e 18 a 21 de 24 verificações manuais atendidas (M10 Regular) \\
	Q9 & Compreende as informações sem legenda? & Autodescritividade & H9: entre 50\% e 75\% de acertos, com o código de horário de T2 abaixo de 50\% (M11 Regular) \\
	Q10 & O estudante fica satisfeito? & Satisfação (em uso) & H10: SUS médio entre 51 e 68 (M12 Regular) \\
	\multicolumn{4}{l}{\textit{G3 -- Flexibilidade}} \\
	Q11 & As tarefas funcionam nos navegadores mais usados? & Adaptabilidade (software) & H11: 22 ou 23 das 24 combinações sem falha (M13 Bom) \\
	Q12 & O SIGAA se adapta a celular, tablet e toque? & Adaptabilidade (hardware) & H12: até 5 de 12 verificações adaptadas (M14 Insatisfatório) e tempo no celular entre 1,5 e 2,0 vezes o do computador (M15 Regular) \\
	\multicolumn{4}{l}{\textit{G4 -- Eficiência de Desempenho}} \\
	Q13 & As páginas carregam em tempo aceitável, em rede comum e móvel? & Comportamento temporal & H13: LCP até 2,5~s em rede comum (M16 Bom) e acima de 2,5~s em rede móvel simulada (M17) \\
	Q14 & O uso de recursos é compatível com rede móvel e plano de dados limitado? & Utilização de recursos & H14: até 1~MB por página (M18 Bom), mas mais de 50 requisições (M19 Regular) \\
\end{longtable}
}
\vspace{-0.6cm}
\begin{center}\small Fonte: Elaborado pela equipe.\end{center}

\section{Métricas}
\label{sec:metricas}

As métricas foram escolhidas por serem simples (contagens, razões ou medianas), objetivas (dois avaliadores seguindo o mesmo procedimento obtêm o mesmo valor) e válidas (cada uma mede o atributo perguntado). Sempre que possível, seguem as medidas da ISO/IEC 25023 e da ISO/IEC 25022 \cite{iso25023,iso25022}. Antes da coleta, a equipe fixa o caminho mais curto de cada tarefa, a definição de passo (clique ou toque que muda o estado da interface), de ajuda (qualquer intervenção do moderador, de colega ou de fonte externa) e de tarefa concluída (resultado correto, sem ajuda, em até 5 minutos). Toda medição registra data, horário, dispositivo, navegador e rede. A \autoref{tab:metricas} resume as 20 métricas.

{\scriptsize
\begin{longtable}{L{0.6cm} L{3.6cm} L{4.5cm} L{5.6cm}}
	\caption{Métricas da avaliação}
	\label{tab:metricas} \\
	\toprule
	\textbf{ID} & \textbf{Métrica (questão)} & \textbf{Cálculo} & \textbf{Fonte e coleta} \\
	\midrule
	\endfirsthead
	\toprule
	\textbf{ID} & \textbf{Métrica (questão)} & \textbf{Cálculo} & \textbf{Fonte e coleta} \\
	\midrule
	\endhead
	\bottomrule
	\endfoot
	M1 & Cobertura funcional (Q1) & $1 - A/B$; $A$ = funções ausentes, $B$ = funções da lista de referência & Lista montada do questionário, das ferramentas comunitárias e das tarefas; verificada por dois avaliadores \\
	M2 & Correção funcional (Q2) & $1 - A/B$; $A$ = verificações divergentes & Comparação com o projeto pedagógico, o calendário oficial e o histórico emitido; dois avaliadores \\
	M3 & Uso de ferramentas externas (Q3) & $A/B$; $A$ = respondentes que usam ferramenta externa & Questionário anônimo (\autoref{sec:questionario}) \\
	M4 & Acerto na primeira escolha de menu (Q4) & $A/B$; $B$ = tentativas & Gravações do teste com usuários \\
	M5 & Conclusão de tarefas sem ajuda (Q5) & $A/B$; $A$ = tarefas concluídas & Gravações e formulário do moderador; também por perfil \\
	M6 & Tempo relativo de execução (Q6) & $T_p / T_e$ & Gravações; $T_e$ = média de três execuções de um especialista \\
	M7 & Desvio de caminho (Q6) & $C_p / C_{min}$ (cliques) & Gravações; $C_{min}$ fixado antes do teste \\
	M8 & Recuperação de erros (Q7) & $A/B$; $A$ = erros recuperados sem ajuda nem retrabalho & Gravações do teste com usuários \\
	M20 & Cenários de erro tratados (Q7) & $A/6$ & Inspeção de seis cenários (expiração de sessão, Voltar, recarregar, nova aba, função fora do período, bloqueio de PDF) \\
	M9 & Erros automáticos de acessibilidade (Q8) & erros / páginas & Extensão WAVE nas páginas finais de T1 a T6 \\
	M10 & Critérios manuais de acessibilidade (Q8) & $A/24$ & Teclado, foco, zoom de 200\% e leitor de tela em seis páginas; dois avaliadores \\
	M11 & Interpretação sem legenda (Q9) & $A/B$; três perguntas por participante & Perguntas de compreensão em T2, T4 e T6 \\
	M12 & Pontuação SUS (Q10) & escala SUS (0 a 100) & Questionário SUS após o teste \cite{brooke1996} \\
	M13 & Compatibilidade entre navegadores (Q11) & $A/24$ (6 tarefas $\times$ 4 navegadores) & Execução pela equipe, com captura de tela \\
	M14 & Páginas adaptadas a telas pequenas e toque (Q12) & $A/12$ (6 páginas $\times$ 2 larguras) & Celular (360~px) e tablet (768~px): rolagem, fonte, sobreposição, toque e alvo de 24~px \\
	M15 & Tempo relativo no celular (Q12) & $T_{celular} / T_{computador}$ & Execução cronometrada pelos avaliadores \\
	M16 & LCP em rede comum (Q13) & média das medianas (s) & Lighthouse, modo \textit{Desktop}, cinco execuções por página \cite{lighthouse} \\
	M17 & LCP em rede móvel simulada (Q13) & média das medianas (s) & Lighthouse, modo \textit{Mobile} \\
	M18 & Peso transferido por página (Q14) & média das medianas (MB) & Lighthouse (\textit{total byte weight}) \\
	M19 & Requisições por página (Q14) & média das medianas & Lighthouse (\textit{network requests}) \\
\end{longtable}
}
\vspace{-0.6cm}
\begin{center}\small Fonte: Elaborado pela equipe, com base em \citeonline{iso25023} e \citeonline{iso25022}.\end{center}

\section{Níveis de pontuação e critérios de julgamento}
\label{sec:niveis}

Todas as métricas usam uma escala ordinal de quatro níveis: Excelente (3), Bom (2), Regular (1) e Insatisfatório (0). Um valor igual ao limite entre dois níveis pertence ao nível superior. Quando há referência externa, ela define os limites (\autoref{sec:medidas-referencia}); nos demais casos, os limites foram definidos pela equipe e justificados na GitPage (\autoref{tab:niveis}). Os limites foram fixados antes da coleta e não serão alterados após a obtenção dos dados.

\begin{table}[htb]
	\caption{Níveis de pontuação por métrica}
	\label{tab:niveis}
	\centering
	\scriptsize
	\begin{tabularx}{\textwidth}{L{3.9cm} C{1.6cm} C{2.1cm} C{2.1cm} Y}
		\toprule
		\textbf{Métrica} & \textbf{Excelente} & \textbf{Bom} & \textbf{Regular} & \textbf{Insatisfatório} \\
		\midrule
		M1 Cobertura funcional & 100\% & 90\% a $<$100\% & 75\% a $<$90\% & $<$75\% \\
		M2 Correção funcional & 100\% & 97\% a $<$100\% & 90\% a $<$97\% & $<$90\% ou divergência em item crítico \\
		M3 Ferramentas externas & $\leq$10\% & $>$10\% a 25\% & $>$25\% a 50\% & $>$50\% \\
		M4 Primeira escolha & $\geq$90\% & 75\% a $<$90\% & 50\% a $<$75\% & $<$50\% \\
		M5 Conclusão sem ajuda & $\geq$90\% & 78\% a $<$90\% & 60\% a $<$78\% & $<$60\% \\
		M6 Tempo relativo & $\leq$1,5 & $>$1,5 a 2,5 & $>$2,5 a 4,0 & $>$4,0 \\
		M7 Desvio de caminho & $\leq$1,11 & $>$1,11 a 1,33 & $>$1,33 a 2,00 & $>$2,00 \\
		M8 Recuperação de erros & $\geq$90\% & 70\% a $<$90\% & 50\% a $<$70\% & $<$50\% \\
		M20 Cenários de erro & 6 & 5 & 3 ou 4 & 0 a 2 \\
		M9 Erros automáticos/página & 0 & $>$0 a 2 & $>$2 a 5 & $>$5 \\
		M10 Critérios manuais & 24 & 22 ou 23 & 18 a 21 & $\leq$17 \\
		M11 Interpretação sem legenda & $\geq$90\% & 75\% a $<$90\% & 50\% a $<$75\% & $<$50\% \\
		M12 SUS & $\geq$80,3 & 68 a $<$80,3 & 51 a $<$68 & $<$51 \\
		M13 Navegadores & 24 & 22 ou 23 & 18 a 21 & $\leq$17 \\
		M14 Páginas adaptadas & 12 & 10 ou 11 & 6 a 9 & $\leq$5 \\
		M15 Tempo no celular & $\leq$1,2 & $>$1,2 a 1,5 & $>$1,5 a 2,0 & $>$2,0 \\
		M16 e M17 LCP (s) & $\leq$1,0 & $>$1,0 a 2,5 & $>$2,5 a 4,0 & $>$4,0 \\
		M18 Peso (MB) & $\leq$0,5 & $>$0,5 a 1,0 & $>$1,0 a 2,0 & $>$2,0 \\
		M19 Requisições & $\leq$25 & $>$25 a 50 & $>$50 a 80 & $>$80 \\
		\bottomrule
	\end{tabularx}
	\fonte{Elaborado pela equipe, com base em \citeonline{sauro2011}, \citeonline{sauro2016}, \citeonline{wcag22}, \citeonline{webvitals} e \citeonline{nielsen1993}.}
\end{table}

O julgamento é feito do nível quantitativo para o conceitual, em três etapas:

\begin{enumerate}
	\item \textbf{Questão}: recebe o \textbf{menor} nível entre as suas métricas. A regra é conservadora: um atributo só é Bom se todas as medidas que o descrevem forem ao menos Boas.
	\item \textbf{Objetivo}: recebe a \textbf{média} dos níveis das suas questões, para que subcaracterísticas com mais métricas não pesem mais que as outras.
	\item \textbf{Característica}: \textit{atende plenamente} se a média for $\geq 2{,}5$ sem questão abaixo de Bom; \textit{atende} se a média for $\geq 2{,}0$ sem questão Insatisfatória; \textit{atende parcialmente} se a média for $\geq 1{,}0$; \textit{não atende} se a média for menor que 1,0; e \textit{inconclusivo} se mais da metade das questões ficar sem dados. Uma questão Insatisfatória impede o julgamento ``atende'', porque um problema grave não é compensado por bons resultados nas demais.
\end{enumerate}

Regras complementares tratam de casos específicos. Métricas sem dados suficientes ficam fora do cálculo; divergência em menção, IRA ou carga horária pendente leva M2 a Insatisfatório; recortes por perfil, página ou tarefa localizam os problemas, mas não alteram o nível da questão, e recortes com menos de três participantes são apenas descritivos. Para apoiar a STI, cada tarefa recebe um índice igual à média dos níveis de M5, M6, M7 e M1 naquela tarefa, e as tarefas são ordenadas da pior para a melhor. Por fim, \textbf{não é calculada uma nota única} para o SIGAA: uma média geral esconderia os problemas que precisam ser priorizados. O resultado final é o quadro com o julgamento de cada característica, na ordem de prioridade da Fase 1, acompanhado do ranking das tarefas.

\section{Hierarquia GQM}
\label{sec:hierarquia}

A \autoref{fig:gqm} apresenta a hierarquia completa, com os três níveis do GQM. Toda métrica está ligada a uma questão, e toda questão, a um objetivo e a uma hipótese.

\begin{figure}[htbp]
	\centering
	\caption{Hierarquia GQM da avaliação}
	\label{fig:gqm}
	\begin{tikzpicture}[
		y=-0.68cm,
		g/.style={draw=blue!60!black, fill=blue!12, rounded corners, align=center, text width=2.6cm, font=\scriptsize, inner sep=3pt},
		q/.style={draw=teal!60!black, fill=teal!8, rounded corners, align=left, text width=5.4cm, font=\scriptsize, inner sep=2.5pt},
		m/.style={draw=orange!70!black, fill=orange!8, rounded corners, align=left, text width=5.6cm, font=\scriptsize, inner sep=2.5pt},
		lig/.style={-{Stealth[length=1.6mm]}, gray!70!black}]

		% Cabeçalhos dos níveis
		\node[font=\footnotesize\bfseries, align=center] at (0,-1.4) {Nível conceitual\\(objetivos)};
		\node[font=\footnotesize\bfseries, align=center] at (4.7,-1.4) {Nível operacional\\(questões)};
		\node[font=\footnotesize\bfseries, align=center] at (10.9,-1.4) {Nível quantitativo\\(métricas)};

		% Métricas
		\foreach \id/\txt/\y in {
			M1/Cobertura funcional/0, M2/Correção funcional/1, M3/Uso de ferramentas externas/2,
			M4/Acerto na primeira escolha/3, M5/Conclusão sem ajuda/4, M6/Tempo relativo de execução/5,
			M7/Desvio de caminho/6, M8/Recuperação de erros/7, M20/Cenários de erro tratados/8,
			M9/Erros automáticos de acessib./9, M10/Critérios manuais de acessib./10,
			M11/Interpretação sem legenda/11, M12/Pontuação SUS/12, M13/Compatibilidade de navegadores/13,
			M14/Páginas adaptadas/14, M15/Tempo relativo no celular/15, M16/LCP em rede comum/16,
			M17/LCP em rede móvel simulada/17, M18/Peso por página/18, M19/Requisições por página/19}
			\node[m] (\id) at (10.9,\y) {\textbf{\id} -- \txt};

		% Questões
		\foreach \id/\txt/\y in {
			Q1/Completude funcional/0, Q2/Correção funcional/1, Q3/Pertinência funcional/2,
			Q4/Reconhecimento de adequação/3, Q5/Apreensibilidade/4, Q6/Operabilidade/5.5,
			Q7/Proteção contra erros/7.5, Q8/Inclusividade/9.5, Q9/Autodescritividade/11,
			Q10/Satisfação (qualidade em uso)/12, Q11/Adaptabilidade (navegadores)/13,
			Q12/Adaptabilidade (telas e toque)/14.5, Q13/Comportamento temporal/16.5,
			Q14/Utilização de recursos/18.5}
			\node[q] (\id) at (4.7,\y) {\textbf{\id} -- \txt};

		% Objetivos
		\node[g] (G1) at (0,1)     {\textbf{G1}\\Adequação Funcional};
		\node[g] (G2) at (0,7.5)   {\textbf{G2}\\Capacidade de Interação (Usabilidade)};
		\node[g] (G3) at (0,13.75) {\textbf{G3}\\Flexibilidade (Portabilidade)};
		\node[g] (G4) at (0,17.5)  {\textbf{G4}\\Eficiência de Desempenho};

		% Ligações objetivo -> questão
		\foreach \gg/\qq in {G1/Q1, G1/Q2, G1/Q3, G2/Q4, G2/Q5, G2/Q6, G2/Q7, G2/Q8, G2/Q9, G2/Q10,
			G3/Q11, G3/Q12, G4/Q13, G4/Q14}
			\draw[lig] (\gg.east) -- (\qq.west);

		% Ligações questão -> métrica
		\foreach \qq/\mm in {Q1/M1, Q2/M2, Q3/M3, Q4/M4, Q5/M5, Q6/M6, Q6/M7, Q7/M8, Q7/M20, Q8/M9,
			Q8/M10, Q9/M11, Q10/M12, Q11/M13, Q12/M14, Q12/M15, Q13/M16, Q13/M17, Q14/M18, Q14/M19}
			\draw[lig] (\qq.east) -- (\mm.west);
	\end{tikzpicture}
	\fonte{Elaborado pela equipe.}
\end{figure}

\section{Relações entre as características}
\label{sec:relacoes}

O roteiro exige discutir, ao final do projeto, as relações entre as características \cite{ramos2026roteiro}. As relações esperadas na Fase 1 foram convertidas em quatro hipóteses de relação, que não recebem nível de pontuação e serão declaradas estabelecidas ou não estabelecidas na Fase 4 (\autoref{tab:relacoes}).

\begin{table}[htb]
	\caption{Hipóteses de relação entre as características}
	\label{tab:relacoes}
	\centering
	\scriptsize
	\begin{tabularx}{\textwidth}{C{0.8cm} L{3.4cm} Y L{4.2cm}}
		\toprule
		\textbf{ID} & \textbf{Relação} & \textbf{Hipótese e verificação} & \textbf{Relação não estabelecida se} \\
		\midrule
		HR1 & Flexibilidade $\rightarrow$ Capacidade de Interação & O esforço adicional no celular se concentra nas páginas não adaptadas: M15 por tarefa cruzado com M14 & M15 $\leq$ 1,5 ou as tarefas de maior M15 não coincidem com as páginas não adaptadas \\
		HR2 & Eficiência de Desempenho $\rightarrow$ Capacidade de Interação & Tarefas com maior LCP têm maior tempo relativo: ordenação por M16 e por M6 (correlação de Spearman descritiva) & As ordenações divergem \\
		HR3 & Capacidade de Interação $\rightarrow$ Adequação Funcional & Funções existentes, mas difíceis de localizar, são declaradas inexistentes: métrica auxiliar MA1 & Nenhuma ocorrência \\
		HR4 & Eficiência de Desempenho $\rightarrow$ Flexibilidade & O LCP em rede móvel é ao menos o dobro do LCP em rede comum: razão M17/M16 & Razão menor que 2,0 \\
		\bottomrule
	\end{tabularx}
	\fonte{Elaborado pela equipe.}
\end{table}

\section{Questionário com estudantes de diferentes semestres}
\label{sec:questionario}

O teste com usuários mede o comportamento em profundidade, mas com poucos participantes. Para ampliar a amostra e comparar estudantes em diferentes momentos do curso, a especificação inclui um questionário online anônimo, respondido em 5 a 8 minutos, com seis blocos: perfil (semestre, campus e previsão de formatura), acesso (dispositivo, rede e problemas no celular), facilidade percebida em cada tarefa T1 a T6 (escala SEQ de 1 a 5 \cite{sauro2016}), compreensão sem legenda (interpretação de um código de horário e da sigla IRA), satisfação (SUS \cite{brooke1996}) e ferramentas externas. O questionário é a fonte da métrica M3 e de parte da lista de referência de M1 e produz indicadores auxiliares de percepção que, como a MA1, não recebem nível de pontuação e não alteram o julgamento: facilidade percebida por tarefa, acerto na interpretação do horário e da sigla IRA, SUS do questionário, dependência do celular e lentidão percebida. Esses indicadores verificam a hipótese de perfil \textbf{HP1}: estudantes de 1º e 2º semestres relatam menor facilidade nas tarefas e acertam menos a interpretação do código de horário do que estudantes a partir do 5º semestre. A comparação usa medianas e proporções e, se cada grupo tiver ao menos dez respondentes, os testes de Mann-Whitney (facilidade) e exato de Fisher (acerto), com nível de significância de 5\%; caso contrário, é apenas descritiva. O instrumento completo, com o mapeamento de cada item às questões GQM, está na GitPage.
